% ===================================================================
%  TABELLA nr. 2 — Il corp uman. --- La recreaziun. --- Ils gieus.
%  Traduction 2026 d'apres texto/io, controlee sur texto/fr, texto/ca
%  et texto/oc. Consigne : voir texto/rm/10-tabella-01.tex.
%
%  DEUX MOITIES SUR UNE PAGE, ET ELLES NE DISENT PAS LA MEME CHOSE :
%  LE CORPS DIT D'OU VIENT LE ROMANCHE, LES JEUX DISENT OU IL VIT.
%
%  Le corps est latin, et souvent plus vieux que celui des voisins.
%  Le cas le plus net est la tete :
%    LE ROMANCHE : « chau », du latin CAPUT ;
%    l'italien :   « testa », et le francais « tête », tous deux de
%                  TESTA, qui en latin est un POT — un mot de
%                  plaisanterie de soldat qui a chasse le mot propre.
%  Le romanche n'a pas suivi la plaisanterie. Meme chose pour le
%  cerveau, « tscharvè », pour la nuque, pour le « culiez » du cou,
%  pour le « dies » du dos : le fonds est latin, direct, et il a
%  garde des mots que les deux grandes langues ont laisses tomber.
%  C'EST LA CONFIRMATION DU TABLEAU 1 SUR LE VOCABULAIRE LE MOINS
%  EMPRUNTABLE QUI SOIT : on emprunte le nom d'une institution, on
%  n'emprunte pas le nom de sa propre tete.
%
%  Les jeux, eux, sont suisses. « chegls » les quilles vient de
%  l'allemand Kegel, la ou l'italien dit « birilli » ; « stelzas »
%  les echasses vient de Stelzen ; « dracun » le cerf-volant suit le
%  Drache allemand et non l'aquilone italien.
%
%  ET LE VELO DE L'ALINEA 31 DIT LE PLUS. Le romanche ecrit
%  « velo » : c'est un mot francais, et l'on croirait donc a un
%  emprunt au voisin roman. Il n'en est rien. L'italien dit
%  « bicicletta », le francais de France dit « vélo » mais n'a pas de
%  frontiere avec les Grisons ; ce qui dit « Velo » autour du
%  romanche, c'est l'ALLEMAND DE SUISSE, la ou l'allemand d'Allemagne
%  dit « Fahrrad ». Le mot est francais, la route est allemande.
%  LA LANGUE D'UN EMPRUNT NE DIT PAS LA ROUTE QU'IL A PRISE — et
%  c'est la premiere fois de la serie que les deux se separent
%  aussi proprement. La colonne luxembourgeoise avait montre le
%  contraire au tableau 13 : un mot allemand pris tel quel y marquait
%  la distance ; ici un mot francais pris tel quel marque la
%  proximite d'un troisieme.
% ===================================================================

%%K t02-apar-1 apar
\VUsaut{4.0mm}
\VUtitre{15.0pt}{60}{TABELLA nr. 2}
\VUsaut{2.0mm}
\VUtitre{11.4pt}{0}{Il corp uman. --- La recreaziun.}
\VUtitre{11.4pt}{0}{Ils gieus.}

%%K t02-tit-0 sub
\VUtitre{13.2pt}{0}{Il corp uman.}
\VUsaut{2.4mm}
\VUcentre{10.2pt}{0}{\textit{(Simpla lecziun da scienzas natiralas.)}}
\VUsaut{3.0mm}

%%K t02-tit-1 sub pos=t02-01-1
\VUpk{10.2pt}{40}{Il chau.}
\VUsaut{3.0mm}

%%K t02-01-1 p
1. --- Las parts principalas dal corp uman èn: il \VUgras{chau}, il \VUgras{tronc}
ed ils \VUgras{members}.

%%K t02-02-1 p
2. --- Il chau cumpiglia duas parts: il \VUgras{crani}, che cuntegna il
\VUgras{tscharvè} e ch'è cuvrì dals \VUgras{chavels}, e la \VUgras{fatscha}.

%%K t02-03-1 p
3. --- Sin noss dissegn vesain nus ni il crani ni il tscharvè; ma nus vesain fitg
cleramain las parts impurtantas da la fatscha.

%%K t02-04-1 p
4. --- Qua è l'emprim il \VUgras{frunt}, che mussa in intellect pussant, in
patratg adina activ. Las \VUgras{templas} èn fitg libras. L'\VUgras{egl} è vivace e
largamain avert. In dens \VUgras{sursigl} e lungas \VUgras{tschilias} al protegian.

%%K t02-05-1 p
5. --- Qua èn plinavant il \VUgras{nas} ed ils \VUgras{narischs}. Il nas ans serva
per udurar e per respirar. Da mintga vart dal nas èn las \VUgras{massellas}, pli
lunsch las \VUgras{ureglias}. La part sutvart da la fatscha consista da la
\VUgras{bucca} e dal \VUgras{mentun}.

%%K t02-06-1 p
6. --- Nus als vesain betg fitg bain, perquai che las \VUgras{mustatschas} ed ils
\VUgras{favoris} ans als zuppan. Ma nus savain che la bucca ha las duas
\VUgras{leftgas}, la \VUgras{missella sura} e la \VUgras{missella sutta} cun lur
\VUgras{dents}, e la \VUgras{lieunga}. Cun la bucca discurrin e maglian nus. Cun la
lieunga gustain nus las victualias.

%%K t02-07-1 p
7. --- L'egl, l'ureglia, il nas e la bucca èn, sco ils dets, ils organs dals
tschintg senss: il \VUgras{sen da la vista}, il \VUgras{sen da l'udida}, il
\VUgras{sen da l'odur}, il \VUgras{sen dal gust}, il \VUgras{sen dal tast}.

%%K t02-08-1 p
8. --- Il chau è pia ina da las parts las pli interessantas dal corp uman.

%%K t02-tit-2 sub
\VUsaut{4.4mm}
\VUpk{10.2pt}{40}{Il tronc.}
\VUsaut{3.0mm}

%%K t02-09-1 p
9. --- Il \VUgras{culiez} colliescha il chau cun il tronc. El cumpiglia la
\VUgras{gula} e la \VUgras{nuca}.

%%K t02-10-1 p
10. --- Il tronc cumpiglia era blers organs essenzials. En il \VUgras{pèz} èn ils
\VUgras{pulmuns}, cun ils quals nus respirain, ed il \VUgras{cor}, ch'è il center
da la vita. Cur ch'il cor sistescha da batter, mora l'esser vivent.

%%K t02-11-1 p
11. --- Las \VUgras{costas} tegnan il pèz.

%%K t02-12-1 p
12. --- Las autras parts dal tronc èn il \VUgras{flanc}, il \VUgras{dies}, las
\VUgras{spatlas} e l'\VUgras{abdomen} (venter). Nus ans mettain sin il flanc u sin
il dies per durmir, nus purtain las chargias sin nossa spatla.

%%K t02-13-1 p
13. --- En l'abdomen èn il \VUgras{stomi} ed ils \VUgras{intestins}. Els ans
servan per digerir las victualias.

%%K t02-tit-3 sub
\VUsaut{4.4mm}
\VUpk{10.2pt}{40}{Ils members.}
\VUsaut{3.0mm}

%%K t02-14-1 p
14. --- Nus avain quatter \VUgras{members}: dus \VUgras{bratschs} e duas
\VUgras{chommas}. Cun ils bratschs prendain, tegnain, auzain e rimnain nus ils
objects; cun las chommas stain, chaminain e currin nus.

%%K t02-15-1 p
15. --- La \VUgras{spatla} colliescha il bratsch cun il corp. In fitg ferm
\VUgras{muscul}, il biceps, mova il bratsch, ch'il \VUgras{cumbel} colliescha cun
l'\VUgras{avantbratsch}. Il \VUgras{pols} furma la giuntira da la \VUgras{maun},
che finescha cun ils \VUgras{dets} e las \VUgras{unglas}. Sin il maun distinguin
nus ina \VUgras{vena} (ed ina arteria), en la quala fluescha il \VUgras{sang}.

%%K t02-16-1 p
16. --- Las parts da la chomma èn: la \VUgras{cossa}, il \VUgras{schanugl} e la
\VUgras{cava dal schanugl}, la \VUgras{polpa}, da la quala la \VUgras{charn} è
cuvrida da la \VUgras{pel}, il \VUgras{chavigl} ed il \VUgras{pe}. Las
\VUgras{ossa} da la chomma èn fitg grossas e fitg fermas. A la fin dal pe èn ils
\VUgras{dets dal pe}. Las autras parts èn la \VUgras{planta dal pe} ed il
\VUgras{chalcogn}.

%%K t02-17-1 p
17. --- Uschia è la descripziun quasi cumpletta dal corp uman.

%%K t02-17-2 p
\textit{Nota.} --- \textit{Cun agid da l'expressiun: « jau hai mal a... (al chau
etc.) » vegn il magister a pudair reduvrar tut quest vocabulari sut il puntg da
vista dal bun u dal nausch} \VUgras{stadi corporal.}

%%K t02-tit-4 sub
\VUsaut{5.0mm}
\VUpk{10.2pt}{40}{Gimnastica.}
\VUsaut{2.4mm}
\VUcentre{10.2pt}{0}{\textit{(Raquint d'in dals scolars.)}}
\VUsaut{3.0mm}

%%K t02-18-1 p
18. --- Per esser flexibels, agils e da buna sanadad stuain nus exercitar nossas
chommas e noss bratschs. Perquai giogain nus e pratitgain \VUgras{gimnastica}.

%%K t02-19-1 p
19. --- Durant l'emna han lieu quests dus exercizis en la curt da l'institut
(vesair la tabella nr. 1). Ma la gievgia e la dumengia giain nus sin in grond
prà, nua che nus avain plaz per currer e per ans divertir.

%%K t02-20-1 p
20. --- Noss magisters e noss parents ans accumpognan mintgatant là; ma il
\VUgras{magister da gimnastica}\textsuperscript{(12)} dirigia bain ils exercizis.

%%K t02-21-1 p
21. --- Sin il prà, a sanestra da l'entrada, sa chattan ils \VUgras{apparats da
gimnastica}\textsuperscript{(1)}: nus rampignain si la
\VUgras{stgala}\textsuperscript{(2)}, nus ans mettain suvrasut vi dals
\VUgras{rintgs}\textsuperscript{(3)} e vi dal \VUgras{trapez}\textsuperscript{(4)},
nus ans tirain si cun ils mauns fin sin la fin da la \VUgras{stgala da
corda}\textsuperscript{(5)} u da la \VUgras{corda cun nuds}\textsuperscript{(6)}.

%%K t02-22-1 p
22. --- Entant siglian noss cumpogns sur il \VUgras{chaval da
lain}\textsuperscript{(7)} u sur las \VUgras{barras parallelas}\textsuperscript{(11)}.
Auters \VUgras{boxeschan}\textsuperscript{(8)} u
\VUgras{schermeschan}\textsuperscript{(9)} cun mascra e cun il \VUgras{florett}.
Auters, a la fin, auzan \VUgras{halteras}\textsuperscript{(10)}.

%%K t02-tit-5 sub
\VUsaut{4.6mm}
\VUpk{10.2pt}{40}{Ils gieus.}
\VUsaut{3.0mm}

%%K t02-23-1 p
23. --- Enturn ils gimnasts giogan mats e mattas divers \VUgras{gieus}. Las mattas
sa divertan cun lur \VUgras{riev}\textsuperscript{(14)}; ellas siglian cun la
\VUgras{corda}\textsuperscript{(13)} u envidan lur frars e cusrins ad ina partida
da \VUgras{croquet}\textsuperscript{(15)}. Ellas han gudent da bittar davent la
\VUgras{bala} da l'adversari cun lur \VUgras{martè da lain}, gist en il mument
ch'el crai da passar facilmain tras l'\VUgras{artg}!

%%K t02-24-1 p
24. --- Ils mats preferan ils gieus pli musculars. Gugent giogan els
\VUgras{chegls}\textsuperscript{(16)}, ch'els bittan giu abilmain cun ina
\VUgras{bola}\textsuperscript{(17)}. Els na spretschan er betg da giogar cun la
\VUgras{trombla}\textsuperscript{(18)} e cun il \VUgras{bilboquet}\textsuperscript{(19)}
u schizunt il \VUgras{gieu da la rana}\textsuperscript{(20)}. Ma lur gieus
preferids èn bain las \VUgras{billas}\textsuperscript{(21)} e la \VUgras{balla
cunter il mir}\textsuperscript{(22)}, perquai ch'il mir da la
\VUgras{serra}\textsuperscript{(26)} è fitg adattà per far returnar la balla
lediera.

%%K t02-25-1 p
25. --- Il pitschen Gion, che va cun sias \VUgras{stelzas}\textsuperscript{(24)}, è
fitg abil e sa fitg bain tegnair l'equiliber. Sper el giogan intgins cumpogns a
\VUgras{zuppar e tschertgar}\textsuperscript{(25)} en quella serra e davos la
\VUgras{saiv}. Auters preferan l'\VUgras{indivinar dal culp}\textsuperscript{(28)}
u il \VUgras{volan}\textsuperscript{(29)}, che sgola d'ina \VUgras{raketa} a
l'autra.

%%K t02-noto-1 noto
\VUnotes{143.2mm}{%
(*) Ils gieus d'uffants han savens nums puramain naziunals. Nus avain empruvà da
als numnar en ido uschè bain sco pussaivel, sia ans inspirond da l'acziun sezza
che als caracterisescha, sia elegend il num naziunal en quest u en quel pajais,
cur ch'el n'è betg memia nunexact. P. ex.: il \VUgras{gieu dal barigl} (tudestg e
franzos) è il \VUgras{gieu da la rana} \textsuperscript{(20)} (englais), en il
qual ils giugaders empruvan da bittar lur discs regonds tras la bucca da la rana
da metal che stat cun bucca averta sin la chascha (il barigl) pertgirada. Il
\VUgras{sigl sur las spatlas}\textsuperscript{(30)} sa differenziescha dal
\VUgras{sigl sur il dies} mo en quai: per siglir sur il chau, ils giugaders
pesan ils mauns sin las spatlas da lur partenari, entant ch'en il «\VUgras{sigl
sur il dies}» els pesan ils mauns mo sin ses dies. --- U lura èn intgins dals
participants inclinads, entant ch'ils auters siglian sur lur dies pesond ils
mauns sin la spatla; ils emprims ston supportar il colp senza sa plegar, ils
segunds ston siglir senza perder l'equiliber (vesair la tabella sezza).%
}

%%K t02-26-1 p
26. --- En il mez dal prà stan dus pignals. A pe d'in dals dus han set u otg mats
organisà ina partida da \VUgras{sigl sur las spatlas}\textsuperscript{(30)} (*).
Quest gieu n'è betg enconuschent en tut ils pajais; ma tuts san tge ch'il
\VUgras{sigl sur il dies} è, il qual ils mats na giogan betg questa giada.

%%K t02-tit-6 sub
\VUsaut{5.0mm}
\VUpk{10.2pt}{40}{Ils gieus a l'aria libra.}
\VUsaut{3.4mm}

%%K t02-27-1 p
27. Il \VUgras{gieu dal tschorv}\textsuperscript{(31)} è era fitg divertent; mattas
e mats sa mettan a turn la \VUgras{benda} sin ils egls e tschiffan ils malabils
che sa laschan cuprir.

%%K t02-28-1 p
28. --- En il mez dal prà, qua è in gieu fitg franzos ma in pau violent: quai è
l'\VUgras{urs}\textsuperscript{(32)}, bler pli canaschus ch'il gieu uschè quiet dal
\VUgras{dracun}\textsuperscript{(33)} u ch'il gieu englais dal
\VUgras{tennis}\textsuperscript{(34)}.

%%K t02-29-1 p
29. --- Quest ultim sport è la legria dals scolars gronds. Noss magisters sezs al
pratitgan gugent. El pretenda gronda abilitad ed agilitad, ed el ma plascha fitg.
Jau lasch a las mattas ils gieus da la \VUgras{balantschina}\textsuperscript{(35)}.

%%K t02-30-1 p
30. --- Jau n'hai betg gugent in auter gieu englais che pudess daventar privlus.

%%K t02-30-1b p
Jau vi discurrer dal \VUgras{ballape}\textsuperscript{(36)}, che allegrescha mes
frar pli vegl. Jau preferesch noss vegl \VUgras{gieu da las
barras}\textsuperscript{(38)}, tuttina igienic e propi pli pauc brutal.

%%K t02-tit-7 sub
\VUsaut{4.6mm}
\VUpk{10.2pt}{40}{Velo e gieus da balla.}
\VUsaut{3.0mm}

%%K t02-31-1 p
31. --- Er gugent fatsch jau in gir dal prà cun il \VUgras{velo}\textsuperscript{(37)}.

%%K t02-32-1 p
32. --- L'onn che vegn emprend jau da \VUgras{chavaltgar}\textsuperscript{(39)}.
Tge bel ch'in chaval è, che va u trotta cun grazia e che sigla en galopp sur ils
obstachels!

%%K t02-33-1 p
33. --- La stad emprendain nus l'\VUgras{art dal nudar}\textsuperscript{(40)}. Cun
delizia ans tuffain e bagnain nus en l'aua clera e frestga dal flum. Mintgatant
laschain nus a la fin ir in \VUgras{ballun}\textsuperscript{(41)} da palpiri, che
va si majestusamain en l'atmosfera, uschespert ch'el è emplenì cun aria chauda.

%%K t02-34-1 p
34. --- I dat anc blers auters gieus, ma jau na finiss betg da als dumbrar, ed
ins vesa tenor quest curt resumaziun ch'ins na s'algorda betg fitg la gievgia sin
noss bel prà.

%%K t02-34-2 p
N.-B. --- \textit{Il magister vegn a cumplettar facilmain quest chapitel,
descrivend pli detagliadamain mintgin dals gieus ils pli impurtants.}
\VUsaut{6.0mm}
\VUfilet{22mm}
